%
% oracle-protocol.tex
%
% Z Specification for the z-spec Oracle Wire Protocol
% Formal model of the newline-delimited JSON contract between a Lean
% executable oracle and a property-based test harness.
%
% Type-check: fuzz -t oracle-protocol.tex
% Animate:    probcli oracle-protocol.tex -animate 20
%

\documentclass[a4paper,10pt,fleqn]{article}
\usepackage[margin=1in]{geometry}
\usepackage{fuzz}
\usepackage[colorlinks=true,linkcolor=blue,citecolor=blue,urlcolor=blue]{hyperref}

\begin{document}

\title{The Oracle Wire Protocol: A Z Specification}
\author{Formal Model of the Oracle--Harness Contract}
\date{August 2026}
\hypersetup{
  pdftitle={The Oracle Wire Protocol: A Z Specification},
  pdfauthor={Punt Labs},
  pdfsubject={Z Specification},
  pdfcreator={fuzz/probcli}
}
\maketitle

\tableofcontents
\newpage

% ============================================================
\section{Introduction}
% ============================================================

The \texttt{/z-spec:oracle} command generates two programs from a Z
specification. The first is an \emph{oracle}: a Lean executable that
holds the specification's state and answers commands about it. The
second is a \emph{driver}: a property-based test harness that sends the
oracle a random sequence of operations, sends the same sequence to the
implementation under test, and compares the two. The property being
tested is that the implementation accepts exactly what the specification
accepts.

The two programs speak newline-delimited JSON over a pipe. The oracle
emits its state before reading anything; thereafter each command it reads
yields exactly one result line. This document specifies what a result
line is.

\subsection{Why the verdict exists}

A result line carries the state \emph{and} a verdict. The verdict is not
decoration, and the reason it must be there is the subject of this
specification.

An operation whose precondition holds mutates the state. An operation
whose precondition fails leaves the state alone. But an operation whose
precondition holds may also, quite legitimately, leave the state alone
--- withdrawing nothing from an account is a lawful withdrawal of zero.
So the after-state does not determine which of the two happened. If the
line carried the state and nothing else, a rejection and a no-op success
would be byte-identical, and a driver comparing traces could not assert
the one property the harness exists to establish. An implementation that
permitted an illegal withdrawal would produce a trace indistinguishable
from one that correctly refused it, and its tests would pass.

That was a real defect in the prose statement of this protocol, and it
survived review because prose is where it lived.

\subsection{Where the property is kept}

A specification earns nothing by asserting a property beside its schemas;
the assertion is then prose too, in the same register that failed. So
what the wire promises is kept in the schemas, and this section only says
where to look.

\texttt{Line} is the result line: the state, the verdict, and the reason,
declared together as a single object because that is what crosses the
wire. Every operation includes it, so there is one place where the format
of a line is decided. Its predicate

\begin{center}
  $(ok! = zfalse) \iff (reason! \neq \emptyset)$
\end{center}

binds the verdict to the reason, so the two are one bit expressed twice
and neither can drift from the other. The reason is a set of at most one
predicate name rather than an optional field, which is what keeps
\texttt{reason} present on every line --- empty on acceptance, a
singleton on rejection.

That last point is not tidiness. If the reason were declared only on the
rejecting operation, then the mere \emph{presence} of the \texttt{reason}
key would tell an observer the operation had failed, independently of the
verdict. The wire would then carry two discriminators, drivers would
start reading the one they were told not to parse, and the verdict would
become the redundant field --- which is the field a later implementer
deletes.

\subsection{What the $\Xi$ is worth}

The rejecting operation is framed with $\Xi Model$ rather than
$\Delta Model$ plus predicates saying the state does not change. It is
worth being exact about what this buys, because the tempting claim is
false: $\Xi Model$ is \emph{defined} as $\Delta Model$ with
$\theta Model' = \theta Model$, so a predicate elsewhere in the schema
can contradict it. ``The frame makes non-mutation a theorem'' is not
true, and a predicate $balance' = balance + 1$ refutes it.

What $\Xi$ buys is the failure mode. Under $\Delta$ with explicit
equalities, deleting one of those equalities leaves a satisfiable schema
that quietly permits mutation on rejection: it type-checks, it
model-checks, and it is wrong. Under $\Xi$, contradicting the frame makes
the schema unsatisfiable, so the operation becomes disabled --- and a
model checker asked for operation coverage reports it as never fired.
The error stops being invisible and becomes a dead operation. That is a
smaller claim than a theorem and it is the one the notation supports.

% ============================================================
\section{Basic Types}
% ============================================================

\begin{zed}
  [PREDICATE]
\end{zed}

The names of the predicates a specification imposes. A rejection names
the predicate it violated, so that a failure report can quote it.

\texttt{PREDICATE} is deliberately a given set with no structure. The
protocol says a driver must branch on the verdict and must not parse the
reason; giving \texttt{PREDICATE} no operations is that rule stated in
mathematics rather than in a warning.

% ============================================================
\section{Free Types}
% ============================================================

\begin{zed}
  ZBOOL ::= ztrue | zfalse
\end{zed}

The verdict. \texttt{ztrue} is the JSON \texttt{true} of an accepted
operation; \texttt{zfalse} is the \texttt{false} of a rejected one.

% ============================================================
\section{Global Constants}
% ============================================================

\begin{axdef}
  maxBalance : \nat \\
  reserve : \nat \\
  reserveMaintained : PREDICATE
  \where
  maxBalance = 3 \\
  reserve = 1 \\
  reserve \leq maxBalance
\end{axdef}

\texttt{maxBalance} bounds the state and the command inputs so that
probcli can enumerate the reachable states; the protocol itself imposes
no such bound. \texttt{reserve} is the minimum balance the account must
keep. \texttt{reserveMaintained} names the precondition of
\texttt{Withdraw} --- the predicate a rejection reports.

The third constraint, $reserve \leq maxBalance$, is not arithmetic
bookkeeping. \texttt{Init} sets the balance to \texttt{maxBalance} and
the minimum to \texttt{reserve}, so it establishes the state invariant
$minimum \leq balance$ only if the two constants are ordered that way.
Discharging that obligation here, where the constants are chosen, is
cheaper than discovering it as an unreachable initial state.

% ============================================================
\section{State}
% ============================================================

The model under test is the account of the running example in
\texttt{commands/oracle.md}: a balance, and a minimum the balance may not
fall below. Only two things about it matter here. It has a precondition
worth violating, so there is something to reject; and it admits an
operation that is accepted and yet changes nothing, so acceptance and
rejection can share an after-state. Everything else about the model is
incidental to the protocol.

\begin{schema}{Model}
  balance : \nat \\
  minimum : \nat
  \where
  %
  % The invariant the rejecting case exists to protect
  %
  minimum \leq balance \\
  %
  % Bounded for ProB animation
  %
  balance \leq maxBalance
\end{schema}

\begin{schema}{Init}
  Model'
  \where
  balance' = maxBalance \\
  minimum' = reserve
\end{schema}

The oracle emits this state before it reads any command, as an ordinary
result line with the verdict \texttt{ztrue}; that line is
\texttt{ReportState} below.

% ============================================================
\section{The Result Line}
% ============================================================

Everything the oracle writes is a \texttt{Line}. It is declared once,
here, and included by every operation, so no operation can emit anything
else and none can emit less.

\begin{schema}{Line}
  balance! : \nat \\
  minimum! : \nat \\
  ok! : ZBOOL \\
  reason! : \finset PREDICATE
  \where
  %
  % At most one violated predicate is reported
  %
  \# reason! \leq 1 \\
  %
  % The verdict and the reason are one bit, expressed twice
  %
  (ok! = zfalse) \iff (reason! \neq \emptyset)
\end{schema}

The state outputs $balance!$ and $minimum!$ are the \texttt{state} object
of the JSON line. They are outputs rather than the after-state itself
because an observer reads a line, not a schema binding: the wire carries
values, and a specification of the wire has to say which values.

The biconditional is the whole contract in one predicate. Read left to
right it says a rejection must give a reason; right to left, that a
reason may only accompany a rejection. It is also what makes $ok!$
undeletable: strike $ok!$ from the declarations and this predicate no
longer type-checks, so the omission that shipped cannot be reintroduced
by quietly dropping a field.

% ============================================================
\section{Operations}
% ============================================================

\texttt{Withdraw} is the specimen operation. Its precondition is that the
balance left afterwards still covers the minimum. The two schemas below
are the two things the oracle can do with it, and between them they cover
every input: a command is either accepted or rejected, never neither and
never both.

\subsection{AcceptWithdraw}

The precondition holds. The state changes, the line reports the state
that now holds, and the verdict is \texttt{ztrue} with no reason.

The interesting instance is $amount? = 0$. The precondition reduces to
the state invariant and therefore always holds, so a withdrawal of
nothing is always accepted --- and leaves the balance exactly where it
was.

\begin{schema}{AcceptWithdraw}
  \Delta Model \\
  Line \\
  amount? : \nat
  \where
  %
  % Bounded input for ProB animation
  %
  amount? \leq maxBalance \\
  %
  % The precondition of Withdraw
  %
  balance - amount? \geq minimum \\
  %
  balance' = balance - amount? \\
  minimum' = minimum \\
  %
  % The line reports the state after the command
  %
  balance! = balance' \\
  minimum! = minimum' \\
  ok! = ztrue \\
  reason! = \emptyset
\end{schema}

\subsection{RejectWithdraw}

The precondition fails. The line reports the unchanged state, the verdict
is \texttt{zfalse}, and $reason!$ is the singleton naming the predicate
that failed.

\begin{schema}{RejectWithdraw}
  \Xi Model \\
  Line \\
  amount? : \nat
  \where
  %
  % Bounded input for ProB animation
  %
  amount? \leq maxBalance \\
  %
  % The negation of the precondition of Withdraw
  %
  \lnot (balance - amount? \geq minimum) \\
  %
  balance! = balance \\
  minimum! = minimum \\
  ok! = zfalse \\
  reason! = \{ reserveMaintained \}
\end{schema}

\subsection{ReportState}

The startup line: the oracle's state, reported before any command is
read. It is an ordinary accepted line with nothing to report a reason
about.

\begin{schema}{ReportState}
  \Xi Model \\
  Line
  \where
  balance! = balance \\
  minimum! = minimum \\
  ok! = ztrue \\
  reason! = \emptyset
\end{schema}

\subsection{Ambiguity}

The witness. \texttt{Ambiguity} holds in a state when two commands are
available from it --- one accepted that changes nothing, one rejected ---
so that both would report the same state. It is the reason the verdict
exists, written as a schema rather than as a remark.

Because it carries a $\Xi$ frame it is an operation, so a model checker
enumerates it, and its appearance in the coverage report is a
demonstration that the ambiguity is real and reachable. Were the two
cases somehow never to coincide, \texttt{Ambiguity} would never fire and
the coverage report would say so.

\begin{schema}{Ambiguity}
  \Xi Model \\
  accepted? : \nat \\
  rejected? : \nat
  \where
  %
  % Bounded inputs for ProB animation
  %
  accepted? \leq maxBalance \\
  rejected? \leq maxBalance \\
  %
  % accepted? is accepted, and changes nothing
  %
  balance - accepted? \geq minimum \\
  balance - accepted? = balance \\
  %
  % rejected? is rejected
  %
  \lnot (balance - rejected? \geq minimum)
\end{schema}

% ============================================================
\section{The Distinguishability Property}
% ============================================================

The property the protocol exists to guarantee:

\begin{quote}
  An observer reading a result line can always tell a rejection from an
  accepted operation that happened to change nothing.
\end{quote}

It is not asserted here. It is carried by three schemas, and what each
one carries is checkable by a tool rather than by a reader's assent.

\texttt{Ambiguity} says the problem is real. It is satisfiable exactly
where an accepted command that changes nothing and a rejected command
both exist, and it is an operation, so a model checker either fires it or
reports it uncovered. In the state $balance = 1$, $minimum = 1$ --- which
\texttt{Init} reaches by a single withdrawal of $2$ --- it holds with
$accepted? = 0$ and $rejected? = 1$: a withdrawal of nothing is accepted
and leaves $balance = 1$; a withdrawal of $1$ is rejected and leaves
$balance = 1$. Two lines, one state.

\texttt{Line} says the problem is answered. Both lines above include it,
so both carry $ok!$, and its biconditional forces $ok!$ and $reason!$ to
agree. An observer therefore reads one discriminator, present on every
line, and never has to infer a verdict from which keys happen to be
there.

\texttt{RejectWithdraw}'s $\Xi$ frame says the rejection did not lie
about the state it reported.

\subsection{What this does and does not defend against}

The specification is checkable, so it is worth recording what the checks
actually catch --- the failure this document exists to discuss was, after
all, an overclaim about a contract.

Deleting $ok!$ from \texttt{Line} makes the document fail
\texttt{fuzz}, because the biconditional still refers to it. Deleting
$ok!$ from \texttt{Line} \emph{and} from every operation's predicates
still fails, for the same reason. That is the mutation worth defending
against, because it is the one that looks like tidying up.

Deleting the biconditional as well, and $reason!$ with it, yields a
document that type-checks and model-checks and says nothing about
verdicts --- which is the specification that shipped. No amount of Z
prevents that, because a type-checker cannot object to an identifier no
paragraph mentions. A specification enforces what it states; it cannot
enforce that it states it. What has changed is the reading cost of
noticing: the contract is now seven schemas rather than a paragraph in a
document of sixteen hundred lines.

The exhaustiveness of the two cases is a further consequence worth
naming. The rejecting predicate is written $\lnot P$ against the
accepting predicate's $P$, over the same declarations, rather than as the
equivalent $balance - amount? < minimum$. Written that way the argument
is on the page: no input satisfies both and every input satisfies one, so
every command yields exactly one result line. The exception is an input
outside $0 \upto maxBalance$, which neither schema admits --- that is the
animation bound, not a property of the protocol, which accepts any
$amount? : \nat$.

% ============================================================
\section{Correspondence With the Wire}
% ============================================================

The Z above and the JSON of \texttt{commands/oracle.md} are the same
object in two notations. The correspondence:

\begin{center}
\begin{tabular}{ll}
  \textbf{Z} & \textbf{JSON} \\
  \hline
  input variables ($amount?$)   & the \texttt{args} object of a command line \\
  \texttt{Line}                 & one result line \\
  $balance!$, $minimum!$        & the \texttt{state} object of a result line \\
  $ok! = ztrue$                 & \texttt{"ok": true} \\
  $ok! = zfalse$                & \texttt{"ok": false} \\
  $reason! = \emptyset$         & \texttt{"reason": null} \\
  $reason! = \{ p \}$           & \texttt{"reason"} naming $p$ \\
  \texttt{ReportState}          & the startup line, emitted before any command is read \\
\end{tabular}
\end{center}

Every row is a schema or a variable of one, which is the point of
declaring \texttt{Line}: there is no row here that the mathematics does
not produce.

Note the two \texttt{reason} rows. The key is on every line --- null when
the command was accepted --- so that a driver cannot learn the verdict
from the key's presence. That is the contract's rule that \texttt{reason}
is diagnostic, enforced by making the only thing it can tell you the
thing $ok$ already told you.

Serialisation of Z types into JSON values is tabulated in
\texttt{commands/oracle.md}; it is a matter of encoding and carries no
obligations that this specification does not already impose.

Two rules of the wire are outside the model, and are recorded here so
that the omission is deliberate rather than accidental. First,
\emph{one command, one line}: the framing of the pipe is a property of
the transport, not of the operations, and the disjointness argument above
is what makes it well defined. Second, the \emph{content} of a reason is
not modelled: \texttt{PREDICATE} is a given set precisely so that nothing
can be computed from a reason, and a driver that parses one is relying on
something this contract does not offer.

\end{document}
